% GENERATED FILE. Edit canonical lessons, then run npm run generate:books.
% canonical-source-hash: fnv1a64:3b5a45958d29e2c2
% canonical-lessons: FR-C519-romancier, FR-C519-intrigue, FR-C519-denouement, FR-C519-rebondissement, FR-C519-captivant

\chapter{Un bon roman}
\label{ch:fr-un-bon-roman}
\hlchaptermodality{519}

\begin{chapteropening}
\textbf{By the end of this chapter:} \emph{I can describe the plot of a novel.}

Complete the last lesson of chapter 519: I can describe the plot of a novel.
\end{chapteropening}

\section[le romancier]{le romancier — novelist}

\label{lesson:FR-C519-romancier}



\begin{quote}
Before the new one: say the French for work, then the French for to put in a frame.
\end{quote}

\subsection*{You'll want to know: le romancier}
\textbf{le romancier} — \textquotedblleft{}novelist\textquotedblright{}.

\begin{quote}
Ce romancier écrit un livre par an.
\end{quote}

\begin{quote}
\emph{This novelist writes a book a year.}
\end{quote}

A woman is \emph{la romancière}.

\begin{cousinweb}
\emph{le romancier} is built on \emph{roman} \textquotedblleft{}novel\textquotedblright{}, from Old French \emph{romanz}, a story told in the common Romance tongue rather than in Latin.
\end{cousinweb}

\begin{grammarlens}[title={What you've built}]
One more word for telling, explaining and describing at B1.
\end{grammarlens}

\subsection*{Your turn}
\begin{itemize}
\raggedright
  \item \emph{Say it:} \emph{le romancier}
  \item \emph{Say it:} \emph{le romancier}, once more
  \item \emph{Say it:} a sentence of your own with \emph{le romancier}
  \item \emph{Recall:} say the French for work, then the French for to put in a frame, then say \emph{le romancier} again
\end{itemize}

\begin{tcolorbox}[breakable,colback=teal!4,colframe=teal!35!black,title={Before you move on}]
What does le romancier mean? (\textquotedblleft{}Novelist\textquotedblright{}.) Say it once more.
\end{tcolorbox}
\section[l'intrigue]{l'intrigue — plot (of a story); intrigue}

\label{lesson:FR-C519-intrigue}



\begin{quote}
Before the new one: say the French for to put in a frame, then the French for novelist.
\end{quote}

\subsection*{You'll want to know: l'intrigue}
\textbf{l'intrigue} — \textquotedblleft{}plot (of a story); intrigue\textquotedblright{}.

\begin{quote}
L'intrigue du livre se passe dans un petit village.
\end{quote}

\begin{quote}
\emph{The plot of the book takes place in a small village.}
\end{quote}

Feminine. Also a secret scheme.

\begin{cousinweb}
\emph{l'intrigue} was borrowed from Italian \emph{intrigo} (EtymDB), from \emph{intrigare}, from Latin \emph{intrīcāre} \textquotedblleft{}to entangle, to perplex\textquotedblright{} (Lewis \& Short).
\end{cousinweb}

\begin{grammarlens}[title={What you've built}]
One more word for telling, explaining and describing at B1.
\end{grammarlens}

\subsection*{Your turn}
\begin{itemize}
\raggedright
  \item \emph{Say it:} \emph{l'intrigue}
  \item \emph{Say it:} \emph{l'intrigue}, once more
  \item \emph{Say it:} a sentence of your own with \emph{l'intrigue}
  \item \emph{Recall:} say the French for to put in a frame, then the French for novelist, then say \emph{l'intrigue} again
\end{itemize}

\begin{tcolorbox}[breakable,colback=teal!4,colframe=teal!35!black,title={Before you move on}]
What does l'intrigue mean? (\textquotedblleft{}Plot (of a story); intrigue\textquotedblright{}.) Say it once more.
\end{tcolorbox}
\section[le dénouement]{le dénouement — ending, outcome (of a story)}

\label{lesson:FR-C519-denouement}



\begin{quote}
Before the new one: say the French for novelist, then the French for plot.
\end{quote}

\subsection*{You'll want to know: le dénouement}
\textbf{le dénouement} — \textquotedblleft{}ending, outcome (of a story)\textquotedblright{}.

\begin{quote}
Je n'ai pas aimé le dénouement du film.
\end{quote}

\begin{quote}
\emph{I did not like the ending of the film.}
\end{quote}

Literally the untying of the knot.

\begin{cousinweb}
\emph{le dénouement} is built on \emph{dénouer} \textquotedblleft{}to untie\textquotedblright{}, from \emph{dé-} and \emph{nouer} \textquotedblleft{}to tie\textquotedblright{}, inherited from Latin \emph{nōdāre}, from \emph{nōdus} \textquotedblleft{}a knot\textquotedblright{} (Lewis \& Short).
\end{cousinweb}

\begin{grammarlens}[title={What you've built}]
One more word for telling, explaining and describing at B1.
\end{grammarlens}

\subsection*{Your turn}
\begin{itemize}
\raggedright
  \item \emph{Say it:} \emph{le dénouement}
  \item \emph{Say it:} \emph{le dénouement}, once more
  \item \emph{Say it:} a sentence of your own with \emph{le dénouement}
  \item \emph{Recall:} say the French for novelist, then the French for plot, then say \emph{le dénouement} again
\end{itemize}

\begin{tcolorbox}[breakable,colback=teal!4,colframe=teal!35!black,title={Before you move on}]
What does le dénouement mean? (\textquotedblleft{}Ending, outcome (of a story)\textquotedblright{}.) Say it once more.
\end{tcolorbox}
\section[le rebondissement]{le rebondissement — twist, new development}

\label{lesson:FR-C519-rebondissement}



\begin{quote}
Before the new one: say the French for plot, then the French for ending.
\end{quote}

\subsection*{You'll want to know: le rebondissement}
\textbf{le rebondissement} — \textquotedblleft{}twist, new development\textquotedblright{}.

\begin{quote}
Le roman est plein de rebondissements.
\end{quote}

\begin{quote}
\emph{The novel is full of twists.}
\end{quote}

Built on \emph{rebondir} \textquotedblleft{}to bounce back\textquotedblright{}.

\begin{cousinweb}
\emph{le rebondissement} is built on \emph{rebondir} \textquotedblleft{}to bounce back\textquotedblright{}, from \emph{bondir} \textquotedblleft{}to leap\textquotedblright{}, which first meant \textquotedblleft{}to resound\textquotedblright{} and goes back to Vulgar Latin \emph{bombitīre}, from Latin \emph{bombus} \textquotedblleft{}a booming sound\textquotedblright{}.
\end{cousinweb}

\begin{grammarlens}[title={What you've built}]
One more word for telling, explaining and describing at B1.
\end{grammarlens}

\subsection*{Your turn}
\begin{itemize}
\raggedright
  \item \emph{Say it:} \emph{le rebondissement}
  \item \emph{Say it:} \emph{le rebondissement}, once more
  \item \emph{Say it:} a sentence of your own with \emph{le rebondissement}
  \item \emph{Recall:} say the French for plot, then the French for ending, then say \emph{le rebondissement} again
\end{itemize}

\begin{tcolorbox}[breakable,colback=teal!4,colframe=teal!35!black,title={Before you move on}]
What does le rebondissement mean? (\textquotedblleft{}Twist, new development\textquotedblright{}.) Say it once more.
\end{tcolorbox}
\section[captivant]{captivant — gripping, captivating}

\label{lesson:FR-C519-captivant}



\begin{quote}
Before the new one: say the French for ending, then the French for twist.
\end{quote}

\subsection*{You'll want to know: captivant}
\textbf{captivant} — \textquotedblleft{}gripping, captivating\textquotedblright{}.

\begin{quote}
J'ai lu un livre captivant pendant les vacances.
\end{quote}

\begin{quote}
\emph{I read a gripping book during the holidays.}
\end{quote}

Built on \emph{captiver} \textquotedblleft{}to captivate\textquotedblright{}.

\begin{cousinweb}
\emph{captivant} is the present participle of \emph{captiver} \textquotedblleft{}to captivate\textquotedblright{}, from Late Latin \emph{captīvāre} \textquotedblleft{}to take prisoner\textquotedblright{}, built on \emph{captīvus} \textquotedblleft{}a captive\textquotedblright{}, from \emph{capere} \textquotedblleft{}to take\textquotedblright{} (Lewis \& Short).
\end{cousinweb}

\begin{grammarlens}[title={What you've built}]
One more word for telling, explaining and describing at B1.
\end{grammarlens}

\subsection*{Your turn}
\begin{itemize}
\raggedright
  \item \emph{Say it:} \emph{captivant}
  \item \emph{Say it:} \emph{captivant}, once more
  \item \emph{Say it:} a sentence of your own with \emph{captivant}
  \item \emph{Recall:} say the French for ending, then the French for twist, then say \emph{captivant} again
\end{itemize}

\begin{tcolorbox}[breakable,colback=teal!4,colframe=teal!35!black,title={Before you move on}]
What does captivant mean? (\textquotedblleft{}Gripping, captivating\textquotedblright{}.) Say it once more.
\end{tcolorbox}
